\emph{(i)} By Proposition \claimref{std-algebra.witt-basis-car}, $a^{2} = a'^{2} = 0$ and $aa' + a'a = 1$, so $aa'a = a(1 - aa') = a$ and $a'aa' = a'$, $aa'$ and $a'a = 1 - aa'$ are complementary idempotents, and every word in $a, a'$ reduces to one of $1, a, a', aa', a'a$ (a word with two equal neighbours vanishes, and alternating words collapse by $aa'a = a$, $a'aa' = a'$); hence $\Lambda_{W}$ is the stated $\Z$-span, and it is a ring. Put $E_{11} = aa'$, $E_{12} = a$, $E_{21} = a'$, $E_{22} = a'a$: then $E_{ij}E_{kl} = \delta_{jk}E_{il}$ (for instance $E_{12}E_{21} = aa' = E_{11}$, $E_{12}E_{22} = aa'a = a = E_{12}$, $E_{12}E_{12} = a^{2} = 0$), and $E_{11} + E_{22} = 1$, so these are matrix units and $\sum x_{ij}E_{ij} \mapsto (x_{ij})$ is a ring isomorphism of their integer span onto $\Mat_{2}(\Z)$, as in Proposition \claimref{std-algebra.cl11-matrix} over $\R$. The check computes every product of the basis $1, a, a', aa'$ in GA, its integer coordinates, and the multiplicativity of the matrix map.

\emph{(ii)} $e_{+} = a + a'$, $e_{-} = a - a'$ and $e_{+}e_{-} = (a + a')(a - a') = a'a - aa'$, so in the matrix picture $1$, $e_{+}$, $e_{-}$, $e_{+}e_{-}$ are $\begin{psmallmatrix} 1 & 0 \\ 0 & 1 \end{psmallmatrix}$, $\begin{psmallmatrix} 0 & 1 \\ 1 & 0 \end{psmallmatrix}$, $\begin{psmallmatrix} 0 & 1 \\ -1 & 0 \end{psmallmatrix}$, $\begin{psmallmatrix} -1 & 0 \\ 0 & 1 \end{psmallmatrix}$; their integer span is the set of matrices with $\alpha \equiv \delta$ and $\beta \equiv \gamma \bmod 2$, closed under products since $e_{+}^{2} = 1$, $e_{-}^{2} = -1$, $e_{+}e_{-} = -e_{-}e_{+}$, of index $2 \cdot 2 = 4$ in $\Mat_{2}(\Z)$ (the change of basis has determinant $\pm 4$), and containing $2E_{ij}$ for all $i, j$. For the idempotents write $x = \alpha + \beta e_{+} + \gamma e_{-} + \delta e_{+}e_{-}$ with real coefficients. The products $e_{+}e_{-}$, $e_{+}(e_{+}e_{-})$, $e_{-}(e_{+}e_{-})$ all anticommute pairwise with their partners, so $x^{2} = (\alpha^{2} + \beta^{2} - \gamma^{2} + \delta^{2}) + 2\alpha(\beta e_{+} + \gamma e_{-} + \delta e_{+}e_{-})$. If $x^{2} = x$ and the non-scalar part of $x$ is nonzero, comparing it gives $2\alpha = 1$; if it is zero, $x = \alpha \in \{0, 1\}$. Elements of $\Lambda_{O}$ have $\alpha \in \Z$, so their only idempotents are $0$ and $1$. The lattice statements are the products $(xa + ya')^{2} = xy(aa' + a'a) = xy$ and $(xe_{+} + ye_{-})^{2} = x^{2} - y^{2}$: the polar form of $Q$ has Gram matrix $\begin{psmallmatrix} 0 & 1 \\ 1 & 0 \end{psmallmatrix}$ on the basis $a, a'$ and $\begin{psmallmatrix} 2 & 0 \\ 0 & -2 \end{psmallmatrix}$ on $e_{+}, e_{-}$, whose half $\begin{psmallmatrix} 1 & 0 \\ 0 & -1 \end{psmallmatrix}$ is the bilinear form; and $\Z e_{+} + \Z e_{-}$ has index $2$ in $\Z a + \Z a'$ since $e_{\pm} = a \pm a'$ and $2a = e_{+} + e_{-}$. This is the sense in which Proposition \claimref{one-relation.two-presentations}, whose two bases generate the same algebra over any ring containing $\tfrac12$, separates over $\Z$. The check verifies closure, the change-of-basis determinant, $2\Lambda_{W} \subset \Lambda_{O}$, the absence of idempotents other than $0, 1$ among the elements of $\Lambda_{O}$ with coordinates in $\{-3, \dots, 3\}$, and that the ten idempotents of $\Lambda_{W}$ other than $0, 1$ with coordinates in $\{-1, 0, 1\}$ all have scalar part $\tfrac12$.

\emph{(iii)} $nmn = n(1 - nm) = n - n^{2}m = n$ and $mnm = m(1 - mn) = m$; then $(mn)^{2} = m(nmn) = mn$ and $mn + nm = 1$. This is the regularity of Lemma \claimref{transport.regularity-localises-zero-divisors-idempotents} with an explicit generalised inverse. The map $\dagger$ is reversion, an anti-automorphism, composed with the automorphism $e_{-} \mapsto -e_{-}$ induced by the isometry, so it is an anti-automorphism; it sends $a = \tfrac12(e_{+} + e_{-})$ to $\tfrac12(e_{+} - e_{-}) = a'$ and so $(aa')^{\dagger} = a'^{\dagger}a^{\dagger} = aa'$, and likewise $(a'a)^{\dagger} = a'a$; with $aa'a = a$ and $a'aa' = a'$ these are the four Penrose equations. The check verifies the identities for the real null pairs of $\Cl(1,1)$, $\Cl(2,2)$ and $\Cl(1,3)$ and the anti-automorphism property of $\dagger$ on the basis of $\Lambda_{W}$.

\emph{(iv)} $x$ is a scalar, so $x$ and $n$ commute, and $n^{2} = 0$ truncates the binomial expansion: $(x + n)^{k} = x^{k} + kx^{k-1}n$, hence $f(x + n) = f(x) + f'(x)n$. Multiplying on the left by $m$ gives $f'(x)\,mn$. For $g = u/h$ with $h(x) \ne 0$, $h(x + n) = h(x) + h'(x)n$ has inverse $h(x)^{-1} - h'(x)h(x)^{-2}n$, and $u(x + n)h(x + n)^{-1} = g(x) + g'(x)n$ by the quotient rule. When $f$ has integer coefficients and $x \in \Z$, $f(x)$ and $f'(x)$ are integers and $n = a$, $m = a'$, $P = a'a$ lie in $\Lambda_{W}$. The check runs $f = 2x^{3} + 5x$ at five rational points and $g = (x^{3} + 3x)/(x^{2} + 1)$ at three.

\emph{(v)} The radical vector $\rho$ anticommutes with every generator orthogonal to it and squares to $0$, so for a blade $x$ not containing $\rho$, $\rho x \rho = \pm x\rho^{2} = 0$, and for a blade containing $\rho$ the product contains $\rho$ twice and vanishes; by linearity $\rho x \rho = 0$ for all $x$. (In the terms of Proposition \claimref{std-algebra.radical-ideal-and-quotient}, $\rho$ lies in the nilpotent ideal $J$ and $\rho\,\Cl\,\rho = 0$.) So $\rho x \rho = \rho$ has no solution, and $\rho \ne 0$ is not regular. Theorem \claimref{conformal.minimal-nondegenerate-completion} embeds $\Cl(p,q,r)$ in the nondegenerate $\Cl(p{+}r, q{+}r)$, in which $\rho$ spans a hyperbolic plane with a partner; that algebra is regular by (iii) and Lemma \claimref{transport.regularity-localises-zero-divisors-idempotents}. The map $\rho \mapsto a$ from $\Cl(0,0,1)$ to $\Cl(1,1)$ is an algebra map by the universal property, since $a^{2} = 0$, and it carries $f(x + \rho) = f(x) + f'(x)\rho$ to (iv). The check verifies $\rho x \rho = 0$ on every blade for each radical generator of $\Cl(0,0,1)$, $\Cl(1,1,1)$, $\Cl(1,2,1)$, $\Cl(0,0,2)$ and $\Cl(1,3,1)$.
